길이가 $M(M \le 20)$인 이진 문자열 $N(N \le 50\, 000)$개가 주어진다.

주어진 문자열 중 1개 이상을 선택하여, 선택한 모든 문자열을 비트별 배타적 논리합(bitwise XOR)한 결과 문자열이 팰린드롬이 되는 경우의 수를 구하여라.

단, 문자열을 선택하는 순서는 고려하지 않는다. 즉, 고르는 순서만 다르고 선택된 문자열의 구성이 같은 경우는 한 번만 센다. 또한, 입력으로 같은 이진 문자열이 여러 번 주어지더라도, 이들은 서로 구별되는 다른 문자열로 취급한다.

배타적 논리합(bitwise XOR) 또는 팰린드롬에 대한 설명은 하단의 노트를 참고하라.